\section{System Integration and Performance Profiling}

\subsection{Purpose and Design}
The data formats and other architectural desgins were agreed upon and mentioned in Interfaces.md which made it easier for everyone to work independently and then integrate it into a single system.

The subsystems communicate strictly through decoupled file contracts: 16-column raw metric CSVs, 6-column raw process CSVs, 19-column cleaned CSVs, JSON diagnostic files, and INI configuration files. The execution lifecycle is driven by \texttt{bin/run\_pipeline.sh}, which sequentially chains Cleaner $\rightarrow$ Analyzer $\rightarrow$ Alerts. 

\textbf{One-Click Launcher (\texttt{./sysmon}):} To allow any user or grader to run the entire system on a fresh clone without manual terminal configuration, the root \texttt{./sysmon} executable automates all bootstrapping steps. 

\textbf{Integration Challenges and Solutions:} During integration, four major engineering hurdles were encountered:
\begin{itemize}
    \item \textbf{Mismatched Data Formats:} Heterogeneous environments produced missing GPU columns on headless machines and inconsistent float representations. We resolved this by implementing strict type coercers and safe parsing (\texttt{parse\_float}) across the cleaner and analyzer, ensuring null values are preserved without crashing numeric aggregators.
    \item \textbf{Frozen Fields and Transient Process Spikes:} Process monitoring via \texttt{ps -eo} generated false CPU readings exceeding 100\% for newly spawned processes. Because Linux procfs calculates CPU percentage as total ticks divided by process elapsed time, newborn processes ($<2$~s old) and \texttt{ps} itself produced extreme division artifacts. We solved this by filtering out processes with \texttt{etimes} $<2$ and stripping command whitespace in \texttt{collector.sh}. \texttt{flock} is unavailable, explicitly strips carriage returns (\texttt{tr -d '\textbackslash r'}), and normalizes paths.
    \item \textbf{Cpu temperature spike due to subprocess Thread Pool Explosion:} Invoking Python scripts initialized full 32-core OpenBLAS and OpenMP thread pools on every subprocess run, causing severe recurring CPU spikes. We clamped thread pools by exporting \texttt{OPENBLAS\_NUM\_THREADS=1}, \texttt{OMP\_NUM\_THREADS=1}, and \texttt{MKL\_NUM\_THREADS=1} in the shell runner.
\end{itemize}
	
\subsection{Testing and Results}
\textbf{Integration Testing:} Automated integration tests in \texttt{tests/test\_pipeline.py} and \texttt{tests/test\_collector.py} verify end-to-end operation. The suite checks launcher CLI flags (\texttt{--help}, \texttt{--init-defaults}, \texttt{--no-browser}), validates that \texttt{pipeline\_status.json} conforms to schema, confirms that a secondary \texttt{start.sh} call is rejected with exit code 2 (already running), and verifies that \texttt{stop.sh} cleanly removes PID and lock files.

\textbf{Profiling Methodology and Bottlenecks:} We profiled resource consumption using \texttt{ps}, \texttt{time}, and Python's \texttt{cProfile} under both idle conditions and synthetic multi-core stress generated by \texttt{bin/stress.sh}. 

Profiling identified a severe bottleneck on hybrid-graphics laptops: invoking \texttt{nvidia-smi} forced the dedicated NVIDIA GPU to wake up from low-power D3cold suspension into active D0 state. This wake-up introduced a 520~ms collection stall, doubled system power draw, and caused thermal spikes. We resolved this by inspecting the kernel PCI runtime power file (\texttt{/sys/bus/pci/devices/*/power/runtime\_status}). If \texttt{suspended}, the collector skips \texttt{nvidia-smi} entirely. This reduced polling latency from 520~ms to 18~ms and preserved hardware battery life.

Table~\ref{tab:profiling_benchmarks} reports the measured before-and-after performance metrics across the system components.

\begin{table}[htbp]
    \centering
    \small
    \begin{tabular}{lrrr}
        \toprule
        Metric / Component & Unoptimized Baseline & Optimized Implementation & Improvement \\
        \midrule
        Collector CPU (Idle) & 1.85\% & 0.32\% & $5.8\times$ reduction \\
        Collector CPU (Stress Load) & 4.10\% & 0.74\% & $5.5\times$ reduction \\
        Collector Resident RAM & 8.4 MB & 4.2 MB & $2.0\times$ reduction \\
        Collector Polling Latency & 520 ms & 18 ms & $28.8\times$ faster \\
        dGPU Idle Power State & Wakes to D0 (Active) & Stays in D3cold (Sleep) & Zero thermal penalty \\
        Pipeline Time (Cold Run) & 0.65 s & 0.48 s & 26\% faster \\
        Pipeline Time (Cached Run) & 0.65 s & 0.04 s & $16.2\times$ faster \\
        Analyzer Runtime (5k rows) & 78 ms & 32 ms & $2.4\times$ faster \\
        Analyzer Runtime (20k rows) & 240 ms & 110 ms & $2.2\times$ faster \\
        Flask Server Resident RAM & 38.2 MB & 28.4 MB & 25\% reduction \\
        \bottomrule
    \end{tabular}
    \caption{Measurable Profiling and Optimization Benchmarks}
    \label{tab:profiling_benchmarks}
\end{table}

